\section*{Chapter \ref{ch:g11:igneous-rocks} --- Igneous Rocks}

\begin{solution}{exo:g11:igneous-rocks:1}
A \omterm{def:g11:igneous-rocks:volcanic-plutonic}{volcanic rock} formed from a magma that solidified at the surface, as a
lava; a \omterm{def:g11:igneous-rocks:volcanic-plutonic}{plutonic rock} crystallised at depth, inside other rocks. Pairs: basalt
and \omterm{def:g11:ocean-ridges:gabbro}{gabbro} (\omterm{def:g11:igneous-rocks:felsic-mafic}{mafic}), \omterm{def:g11:igneous-rocks:rock-names}{andesite} and \omterm{def:g11:igneous-rocks:rock-names}{diorite} (\omterm{def:g11:igneous-rocks:felsic-mafic}{intermediate}), \omterm{def:g11:igneous-rocks:rock-names}{rhyolite} and granite
(\omterm{def:g11:igneous-rocks:felsic-mafic}{felsic}).
\end{solution}

\begin{solution}{exo:g11:igneous-rocks:2}
(1) Glassy. (2) Coarse-grained. (3) \omterm{def:g11:igneous-rocks:porphyritic}{Porphyritic} (olivine \omterm{def:g11:igneous-rocks:porphyritic}{phenocrysts} in a
fine groundmass). (4) Fine-grained.
\end{solution}

\begin{solution}{exo:g11:igneous-rocks:3}
\qty{48}{\percent}: \omterm{def:g11:igneous-rocks:felsic-mafic}{mafic}, \omterm{def:g11:ocean-ridges:gabbro}{gabbro}. \qty{60}{\percent}: \omterm{def:g11:igneous-rocks:felsic-mafic}{intermediate}, \omterm{def:g11:igneous-rocks:rock-names}{diorite}.
\qty{72}{\percent}: \omterm{def:g11:igneous-rocks:felsic-mafic}{felsic}, granite.
\end{solution}

\begin{solution}{exo:g11:igneous-rocks:4}
A \omterm{def:g11:igneous-rocks:intrusion}{dyke} is a sheet of igneous rock that cuts across the layers of the rocks it
intrudes; a \omterm{def:g11:igneous-rocks:intrusion}{sill} is a sheet parallel to the layers, between them; a \omterm{def:g11:igneous-rocks:intrusion}{pluton} is
a large \omterm{def:g11:igneous-rocks:intrusion}{intrusion} of any shape formed at depth. The vertical \omterm{def:g11:igneous-rocks:intrusion}{dyke} that rises
from the \omterm{def:g11:igneous-rocks:intrusion}{pluton} fed the volcano.
\end{solution}

\begin{solution}{exo:g11:igneous-rocks:5}
Olivine forms first; potassium feldspar, quartz and sometimes muscovite
last. Plagioclase starts rich in calcium and becomes richer and richer in
sodium as the magma cools.
\end{solution}

\begin{solution}{exo:g11:igneous-rocks:6}
(1) A: \qty{74}{\percent} \ce{SiO2}, \omterm{def:g11:igneous-rocks:felsic-mafic}{felsic}; B: \qty{47}{\percent}, \omterm{def:g11:igneous-rocks:felsic-mafic}{mafic}.
(2) A is poor in \ce{FeO}, \ce{MgO} and \ce{CaO} and rich in \ce{Na2O} and
\ce{K2O}, as a \omterm{def:g11:igneous-rocks:felsic-mafic}{felsic} magma is; B is rich in \ce{FeO}, \ce{MgO} and \ce{CaO},
poor in \ce{Na2O}, as a \omterm{def:g11:igneous-rocks:felsic-mafic}{mafic} magma is. (3) A is a \omterm{def:g11:igneous-rocks:rock-names}{rhyolite}; B a \omterm{def:g11:ocean-ridges:gabbro}{gabbro}.
\end{solution}

\begin{solution}{exo:g11:igneous-rocks:7}
The time grows as the square of the thickness. \qty{1}{m}: $10/4 = 2.5$
days. \qty{20}{m}: $10 \times 10^2 = 1000$ days, about 2.7 years.
\qty{200}{m}: $10 \times 100^2 = \num{100000}$ days, about 270 years.
\end{solution}

\begin{solution}{exo:g11:igneous-rocks:8}
(1) \unit{m^2} divided by \unit{m^2/s} gives seconds. (2) $400/10^{-6} =
\qty{4e8}{s}$, that is $4 \times 10^8 / 3.16 \times 10^7 \approx 12.7$
years. (3) $0.21 \times 12.7 \approx 2.7$ years.
\end{solution}

\begin{solution}{exo:g11:igneous-rocks:9}
It is \omterm{def:g11:igneous-rocks:porphyritic}{porphyritic}: plagioclase \omterm{def:g11:igneous-rocks:porphyritic}{phenocrysts} in a fine-grained groundmass.
First the magma cooled slowly at depth, and the plagioclase crystals grew
large while most of the magma stayed liquid; then it rose and cooled fast,
and the remaining liquid crystallised as the fine groundmass.
\end{solution}

\begin{solution}{exo:g11:igneous-rocks:10}
(1) Left: 2 Mg, 2 Si, 6 O; right: 2 Mg, 2 Si, 6 O. (2) \qty{1}{kg} of
olivine is $1000/140.7 \approx 7.11$~mol, which takes as many moles of
silica: $7.11 \times 60.1 \approx \qty{427}{g}$. (3) With no silica left in
the liquid, the reaction cannot go on: the olivine stays.
\end{solution}

\begin{solution}{exo:g11:igneous-rocks:11}
(1) $2.5 \times 20 \times 1100 = \qty{55000}{MJ}$, \qty{5.5e10}{J}.
(2) $900 \times 20 = \qty{18000}{MJ}$, \qty{1.8e10}{J}. (3) The heat of
crystallisation is about a third of the heat lost by cooling, a quarter of
the total: enough to slow the freezing a great deal, because it is all
released above the solidus.
\end{solution}

\begin{solution}{exo:g11:igneous-rocks:12}
(1) The smallest crystals at the contact, which froze within a day; the
largest at the centre, which stayed liquid for about 950 days. (2) No: the
first \qty{5}{m} took about 545 days, the last \qty{5}{m} about 400 days; the
front slows down at first, then speeds up near the centre, where the heat
leaving through both sides meets. (3) Even at its centre, the \omterm{def:g11:igneous-rocks:intrusion}{sill} froze in
under three years; a granite with large crystals needs a far longer time.
\end{solution}

\begin{solution}{exo:g11:igneous-rocks:13}
(1) The \omterm{def:g11:igneous-rocks:intrusion}{pluton} is 100 times thicker: $2.6 \times 100^2 = \num{26000}$
years. (2) For example: a deep \omterm{def:g11:igneous-rocks:intrusion}{pluton} sits in rock that is already hot, so
it loses heat more slowly than into rock at \qty{50}{\degreeCelsius}; it may
be fed by new magma again and again; a \omterm{def:g11:igneous-rocks:intrusion}{pluton} is not an infinite sheet;
the properties of rock change with temperature.
\end{solution}

\begin{solution}{exo:g11:igneous-rocks:14}
A \omterm{def:g11:igneous-rocks:felsic-mafic}{mafic} magma is hot and poor in silica: it crystallises the first minerals
of the series --- olivine, pyroxene, calcium-rich plagioclase --- and runs out
of silica before quartz can form. A \omterm{def:g11:igneous-rocks:felsic-mafic}{felsic} magma is cooler and rich in
silica: olivine is not stable in it, and it goes on to the end of the
series, forming sodium-rich plagioclase, potassium feldspar and quartz.
\end{solution}

\begin{solution}{exo:g11:igneous-rocks:15}
(1) $b = \sqrt{t\kappa/C} = \sqrt{86400 \times 10^{-6}/0.21} \approx
\qty{0.64}{m}$. (2) The law gives $0.21 \times 25/10^{-6} \approx
\qty{5.3e6}{s}$, about two months: the same order as a month (a lava flow
also loses heat to the air at its top). (3) Glassy for material erupted into
air or water; fine-grained for lava flows and thin sheets; coarser, but still
small, for thick shallow \omterm{def:g11:igneous-rocks:intrusion}{intrusions}.
\end{solution}

\begin{solution}{pb:g11:igneous-rocks:1}
\textbf{1.} A \omterm{def:g11:igneous-rocks:intrusion}{sill}: it lies parallel to the beds, between them.

\textbf{2.} At the contact, very fine-grained, almost glassy; at the centre,
fine- to medium-grained, the minerals just visible with a lens.

\textbf{3.} The contact lost its heat fastest, into the cold sandstone: the
crystals had no time to grow.

\textbf{4.} A chilled margin. It proves that the magma was hot when it came
into place, and the sandstone cold.

\textbf{5.} A \omterm{def:g11:igneous-rocks:intrusion}{dyke}, the feeder of the \omterm{def:g11:igneous-rocks:intrusion}{sill}.

\textbf{6.} \omterm{def:g11:igneous-rocks:felsic-mafic}{Mafic}: \qty{49}{\percent} \ce{SiO2}, rich in iron, magnesium and
calcium.

\textbf{7.} Olivine, pyroxene and calcium-rich plagioclase.

\textbf{8.} A basalt; the \omterm{def:g11:igneous-rocks:volcanic-plutonic}{plutonic rock} of the same composition is \omterm{def:g11:ocean-ridges:gabbro}{gabbro}.

\textbf{9.} \omterm{def:g11:igneous-rocks:texture}{Texture} depends on the cooling rate, which differs from the
contact to the centre, not on the composition.

\textbf{10.} Olivine first; plagioclase starting at about the same point as
pyroxene.

\textbf{11.} $b^2/\kappa = 400/10^{-6} = \qty{4e8}{s}$.

\textbf{12.} $0.206 \times 4 \times 10^8 \approx \qty{8.2e7}{s}$, that is
$8.2 \times 10^7 / 3.16 \times 10^7 \approx 2.6$ years.

\textbf{13.} $0.206 \times 4/10^{-6} \approx \qty{8.2e5}{s}$, about 9.5
days.

\textbf{14.} $0.07 \times 4 \times 10^8 = \qty{2.8e7}{s}$, about 0.9 year.
Crystallising magma releases heat, which must also be carried away before
the magma becomes solid.

\textbf{15.} About 950 days against 0.6 days: about 1600 times faster.

\textbf{16.} Between the two: a few years, at the long end of a lava flow,
in the months-to-years of shallow \omterm{def:g11:igneous-rocks:intrusion}{intrusions}.

\textbf{17.} A thousand times longer needs a thickness $\sqrt{1000} \approx
32$ times larger: about \qty{630}{m}.

\textbf{18.} Near the surface, a body of magma would erupt, or cool too fast;
a large body that cools for a very long time must be deep, in rock that is
already hot.

\textbf{19.} They have been worn away by erosion.

\textbf{20.} The \omterm{def:g11:igneous-rocks:intrusion}{sill} took about \textbf{2.6 years} to solidify.
\end{solution}
